% !PRINTED 83
\weberpage{83}

We now make the assumption that among the \(\nu\)-rowed determinants of the matrix \gesperrt{at least one} is different from zero, while the \((\nu+1)\)-rowed and consequently also the higher determinants, if such exist, shall all vanish. \(\nu\) may be any number which is not greater than the smaller of the two numbers \(n\) and \(m\) (or, if \(n=m\), does not exceed this common value).

Such a number \(\nu\) can always be found if we exclude the already disposed of, wholly uninteresting case in which all the coefficients \(a_i^{(k)}\) vanish.

Without limiting generality, we may, to simplify the notation, assume that the non-vanishing \(\nu\)-rowed determinant is
\begin{equation}
        A=\begin{vmatrix}
        a_1^{(1)}&a_2^{(1)}&\cdots&a_\nu^{(1)}\\
        a_1^{(2)}&a_2^{(2)}&\cdots&a_\nu^{(2)}\\
        \vdots&\vdots&&\vdots\\
        a_1^{(\nu)}&a_2^{(\nu)}&\cdots&a_\nu^{(\nu)}
        \end{vmatrix}.
\tag{4}
\end{equation}
For it is plainly open to us to arrange the system of equations (1) in any order, and furthermore we can choose the notation of the unknowns \(x\) so that any \(\nu\) of them are the first \(\nu\). We denote the subdeterminants of \(A\), as before, by \(A_i^{(k)}\), where \(i,k\) range from \(1\) to \(\nu\).

\gesperrt{I. If first \(\nu=n\), which presupposes that \(m\) is not smaller than \(n\), then equations (1) have no solution other than that contained in equations (2).}

For let us take the first \(n\) equations of (1):
\begin{equation}
\begin{aligned}
 a_1^{(1)}x_1+a_2^{(1)}x_2+\cdots+a_n^{(1)}x_n&=0,\\
 a_1^{(2)}x_1+a_2^{(2)}x_2+\cdots+a_n^{(2)}x_n&=0,\\
 &\vdots\\
 a_1^{(n)}x_1+a_2^{(n)}x_2+\cdots+a_n^{(n)}x_n&=0,
\end{aligned}
\tag{5}
\end{equation}
and multiply these in succession by \(A_\mu^{(1)},A_\mu^{(2)},\ldots,A_\mu^{(n)}\), where \(\mu\) may be any one of the indices \(1,2,\ldots,n\), and add them. Then, since by \S.\,22, (2) and (6),
\[
 \sum_{1,\nu}^{i} A_\mu^{(i)}a_\lambda^{(i)}=0\quad\hbox{or}\quad=A
\]
%CONT>>
